\documentclass[a4paper,11pt]{article}
\usepackage[utf8]{inputenc}
\usepackage[T1]{fontenc}
\usepackage[french]{babel}
\usepackage{geometry}
\geometry{left=2cm,right=2cm,top=1.5cm,bottom=1.5cm}
\usepackage{hyperref}
\usepackage{enumitem}

\begin{document}

\begin{flushleft}
    \textbf{AISSAOUI Ismail} \\
    Ingénieur en Intelligence Artificielle \& Data Science \\
    Chlef, Algérie \\
    Tél : +213 660 70 77 96 \\
    Email : \href{mailto:ismail.aissaoui.pro@gmail.com}{ismail.aissaoui.pro@gmail.com} \\
    Portfolio : \href{https://ismail-aissaoui.vercel.app}{ismail-aissaoui.vercel.app}
\end{flushleft}

\begin{flushright}
    À l'attention de \\
    \textbf{M. Bart Lamiroy} (\href{mailto:bart.lamiroy@univ-reims.fr}{bart.lamiroy@univ-reims.fr}) \\
    \textbf{M. Frédéric Blanchard} (\href{mailto:frederic.blanchard@univ-reims.fr}{frederic.blanchard@univ-reims.fr}) \\
    \textbf{Mme. Dorsaf Zekri} (\href{mailto:dorsaf.zekri@univ-reims.fr}{dorsaf.zekri@univ-reims.fr}) \\
    Laboratoire CReSTIC, URCA \\
    Fait le \today
\end{flushright}

\vspace{0.5cm}

\noindent \textbf{Objet : Candidature à la thèse « Graphes de patients et représentations latentes pour l'aide au raisonnement clinique »}

\vspace{0.5cm}

Madame, Messieurs,

\medskip

C'est avec un grand enthousiasme que je vous soumets ma candidature pour le projet de thèse portant sur la modélisation de graphes de patients et l'apprentissage de représentations latentes. Actuellement en dernière année de diplôme d'ingénieur en Intelligence Artificielle à l'École Supérieure en Informatique (ESI-SBA), je suis particulièrement motivé par le défi de concevoir des systèmes d'IA transparents, robustes et alignés avec le raisonnement humain.

Votre proposition de recherche résonne profondément avec mes travaux. Fasciné par la structuration des relations sémantiques, j'ai notamment développé une composante métier de « Graph-Based Similarity Intelligence », exploitant les Réseaux de Neurones sur Graphes (GNN) et les modèles de langage (SBERT) pour apprendre des relations complexes entre profils et optimiser l'appariement. Bien qu'appliquée originellement aux ressources humaines, l'approche sous-jacente — construire un graphe où les nœuds sont des entités et les arêtes encodent des similarités apprises — est conceptuellement identique à la modélisation de cohortes de patients.

De plus, mon intérêt pour la dimension « utilisateur » et l'interprétabilité des algorithmes (XAI) m'a conduit à mener des projets sur la détection d'anomalies (Elliptic++) et l'adaptation incrémentale par feedback actif (Recognizini). Ce dernier système génère et met à jour des embeddings latents via un processus incluant l'évaluation d'experts humains, une mécanique très proche du raisonnement clinique fondé sur la comparaison de cas prototypiques et la détection d'atypicités.

L'exigence de rigueur scientifique face aux données multimodales et critiques ne m'est pas étrangère. Au cours de mon stage de fin d'études chez \textbf{Sonatrach}, j'ai été confronté à l'analyse de signaux complexes pour la prédiction de défaillances industrielles à haut risque via des jumeaux numériques. Cette expérience m'a inculqué la nécessité vitale d'une IA de confiance, capable d'être interprétée par un expert métier (ici, l'ingénieur de fiabilité ; dans votre contexte, le médecin).

Le cadre du laboratoire CReSTIC correspond idéalement à mes aspirations de recherche. Ma maîtrise des frameworks (PyTorch Geometric, HuggingFace) et ma rigueur dans l'implémentation de pipelines ML me permettront d'être opérationnel de la formalisation théorique jusqu'au prototypage. 

Je vous prie de trouver ci-joint mon CV détaillé, qui met en évidence la pertinence de mon profil pour vos objectifs. Je reste à votre entière disposition pour un entretien afin de discuter plus amplement de ma vision sur ce projet.

Je vous prie d'agréer, Madame, Messieurs, l'expression de mes salutations respectueuses.

\vspace{0.4cm}

\begin{flushleft}
    \textbf{AISSAOUI Ismail}
\end{flushleft}

\end{document}
